\documentclass[11pt,a4paper]{article}
\usepackage[T1]{fontenc}
\usepackage[utf8]{inputenc}
\usepackage{lmodern}
\usepackage{amsmath,amssymb,amsthm,mathtools,mathrsfs}
\usepackage{graphicx,float,xcolor,multicol}
\usepackage[shortlabels]{enumitem}
\usepackage[margin=27mm]{geometry}
\usepackage{microtype,needspace,placeins}
\usepackage{tikz,tikz-cd}
\usetikzlibrary{arrows.meta,calc,positioning,patterns,decorations.pathreplacing}
\usepackage[colorlinks=true,linkcolor=teal,urlcolor=teal,citecolor=teal]{hyperref}
\setlength{\parindent}{0pt}
\setlength{\parskip}{.6em}
\setcounter{secnumdepth}{0}
\allowdisplaybreaks
\raggedbottom
\providecommand{\cross}{\times}
\DeclareUnicodeCharacter{2020}{\ensuremath{\dagger}}
\DeclareMathOperator{\supp}{supp}
\DeclareMathOperator*{\esssup}{ess\,sup}
\providecommand{\chapter}{\section}
\newenvironment{tool}[2]{\subsection*{Tool #1 --- #2}}{}
\newcommand{\ToolRef}[1]{Tool #1}
\newcommand{\FigureTag}[2]{\textit{Figure #1}}
\newcommand{\Result}[1]{\subsection*{#1}}
\newcommand{\Application}[1]{\subsubsection*{Application\if\relax\detokenize{#1}\relax\else\space--- #1\fi}}
\newcommand{\ProofHeading}{\par\textit{Proof.}\space}
\newcommand{\ProofOf}[1]{\par\textit{Proof of #1.}\space}
\newcommand{\RemarkHeading}{\par\textbf{Remark.}\space}
\newcommand{\NoteHeading}{\par\textbf{Note.}\space}
\newcommand{\ExerciseHeading}{\par\textbf{Exercise.}\space}

\begin{document}
\section*{Scalpel, Forceps, and other Surgical Tools}
% BLOG-CONTENT-BEGIN

Let us look at the aspects of quasiconformal mapping theory that enter holomorphic dynamics. The material is treated in depth in \cite{ahlfors2006lectures}, \cite{branner2014quasiconformal}, \cite{hubbard2016teichmuller}, and \cite{imayoshi2012introduction}.

Dynamics of a given holomorphic self map $f$ on a Riemann surface $S$ is a topological invariant. For instance, fixed points, critical points, the nature of the periodic points - attracting, repelling, superattracting, parabolic, or non-linearizable, all remain the way they were upon one's change of coordinates on $S$ via a homeomorphism. The issue is rigidity. Conformal maps are too rigid, that is, they preserve way too much information. A biholomorphic change of coordinate preserves all the analytic parameters of the dynamics of $f$, like the multipliers, conformal description of Fatou components, holomorphic degree, etc.

To gain some flexibility, we first introduce a way to measure how conformal a map is and then examine what happens when the biholomorphic coordinate change $\phi:S\to S'$ is allowed a controlled amount of distortion. The map $\phi$ must remain a homeomorphism so that the topological core of the dynamics is preserved, while its analytic parameters are allowed to vary. These almost-conformal maps are called quasiconformal maps. We begin by measuring a function's conformality.

The details are developed in \BlogPost{quasiconformal-mappings}{Quasiconformal Mappings}, \BlogPost{the-mapping-theorem-and-invariant-complex-structures}{The Mapping Theorem and Invariant Complex Structures}, \BlogPost{prime-ends-and-landing-of-external-rays}{Prime Ends and Landing of External Rays}, and \BlogPost{boundary-behaviour-of-quasiconformal-maps}{Boundary Behaviour of Quasiconformal Maps}.

Having laid out the theory necessary for surgical constructions let us briefly reiterate it by listing the essential tools one would need. In order, they are:
\begin{enumerate}
    \item \textbf{Measurable Riemann Mapping Theorem}. It shall play the role of a scalpel by instigating any surgical procedure.
    \item \textbf{Invariant Beltrami Forms} will be our forceps, holding the incision for as many deformations as one would desire.
    \item \textbf{Theorems on Boundary and Domain Extensions} shall be the actual tools that one would need to fundamentally change the nature of a self map.
    \item \textbf{Weyl's Lemma} returns the construction to the holomorphic category.
\end{enumerate}

Now, there are certain rules to follow. One should know when these tools can successfully be applied, and that too in the order that they have been listed above. For that, we have the following theorems which shall in some sense be our principles. The theorems illustrate instances where some of these tools can be applied, hence we shall call them surgical principles\footnote{The ones presented here are different from Shishikura's Principles. They are not needed for the constructions that follow.}.

\subsection*{Theorem 3.27}{(Principle I)}: If $\Gamma$ is a quasiarc and $\phi: U \rightarrow V$, a homeomorphism that is $K$-qc on $U\backslash \Gamma$. Then $\phi$ is $K$-qc on $U$. (Note: As a consequence, if the boundary pieces are regular enough, one can construct global qc maps by gluing them along such boundary pieces).

\subsection*{Theorem 3.28}{(Principle II)}: Let $U \subset \mathbb{C}$ be open, $C \subset U$ compact, $\phi$ and $\Phi$ two mappings $U\rightarrow \mathbb{C}$ which are homeomorphisms. Suppose $\phi$ is qc, and so is $\Phi$ on $U\backslash C$, and that $\phi = \Phi$ on $C$. Then $\Phi$ is quasiconformal. (Note: This is referred to as Rickman's Lemma).

\subsection*{Theorem 3.29}{(Principle III)}: Let $f: \hat{\mathbb{C}} \rightarrow \hat{\mathbb{C}}$ be quasiregular. If $\exists$ $K$, such that $f^{\circ n}$ is $K$-quasiregular, $\forall n$, $n \geq 1$. Then $f$ is quasiconformally conjugate to a rational map. (Note: In literature, this is called the Straightening Theorem. It informs us of when the incision done on $f$ can lead us to a situation where Weyl's Lemma can be applied).

\BlogPost{first-surgery}{\emph{Sullivan's No Wandering Theorem}}, \BlogPost{second-surgery}{\emph{The Multiplier Map}}, and \BlogPost{third-surgery}{\emph{McMullen's Surgery}} present three different instances of quasiconformal surgery in complex dynamics.
% BLOG-CONTENT-END
\bibliographystyle{amsalpha}
\bibliography{tools/profiles/surgeries/MSc_Dissertation}
\end{document}
